\section{状态、映射与参考实现}

本节采用\textbf{二进制/SLC PCM、八位平面、电压式近似 ACIM}，配32个实际电流写驱动及复用读侧 ADC 的终态校验。共同28 nm外围下，成对工程情景得到 $\rho\approx1.04$--$5.21\,\mathrm{MB/s}$、$\tau\approx4.95$--$7.97\,\mathrm{MB/s}$、$\RI^*\approx0.211$--0.653；参考点分别为 $2.500\,\mathrm{MB/s}$、$6.874\,\mathrm{MB/s}$ 和0.364。结果来自一份完整局部参考设计的操作预算，既不是原论文芯片的实测两路性能，也不是 PCM 材料固有范围。

原生参考选取一个 $256\,\text{rows}\times1024\,\text{columns}$ sub-bank，尺寸来自\sid{PCM-03} PDF p.3 Fig.11.3.7。八个相邻SLC cell表示一个INT8权重，得到 $W[128,256]$；$B_S=256$ Byte，有效resident容量32768 Byte，共262144个cell。每位平面为256行128列，交付128个24-bit结果容器。符号位按 $-2^7$ 数字重构。二态只约束存储类别，不保证SET cell电导完全相同；本节采用校准近似部分和，不把容器位宽当作物理精度。

\sid{PCM-03} 的40 nm宏以2个SLC和3个MLC cell表示8-bit权重，同次激活8条WL，8-bit输入串行八拍。各列经固定偏置电流 $I_M$ 产生电压 $V_{\rm MAC}$，再用四阶段VSR-VSA及移位加法读出。其纯SLC信号裕量对照支持选取二状态实现，但本文不沿用三只MLC或数据相关跳拍。保留\emph{该电压式组织的八项分组}，八个等尺度SLC平面并行，因此输入分32组；八项二进制部分和为 $0\ldots8$，共有9个等级。8行是所选实现的结构条件，不能升级为 PCM 的介质上限。

写侧采用32个实际IDAC，电流和波形以\sid{PCM-01}为锚点。该源芯片用对角栅线避免任一BL/SL电流拥挤，32个IDAC是其实际面积/功耗选择；本参考不把该对角gate直接嫁接到电压求和阵列。\sid{PCM-06} PDF p.4 Fig.6给出选一行、各列独立编程的机制依据；这里只借行选通原则，不引入其512列资源或模拟闭环算法。

本bank保留原生compute WL：写/verify时只打开一行 $r$，32个IDAC选择逻辑列 $c=4h+s$。当前bit平面 $w$ 的物理列为 $p=8c+w$，只有这些32列获得编程压差；同一行的其余列BL/SL等电位inhibit，其他255行的access gate关闭。每列仅一个目标，不需要每cell增加diagonal gate mux。列外围实际配置读/写隔离开关、八平面路由、耐压选择器和每IDAC四列mux。单行公共回流必须有足够宽金属和受控回流端，承受RESET的 $32\times700\ \mu$A=22.4 mA与SET的4 mA，并保持IDAC合规范围；不能只因目标坐标不同就忽略共线压降。批配置及三段驱动预算包含上述WL/列切换与恢复条件。

八位平面是一份矩阵。读时断开IDAC，由固定 $I_M$、钳位及电压前端服务；写时隔离读前端，verify时同一WL只提供每列一个cell。这个行选通、供给及电压读验组合是明确资源条件的跨实现参考，不称为PCM-01/03同一芯片的实测复原。

\section{完整 streaming 求值与前端预算}

沿用R0的每平面16 ADC、共128 ADC、16条数字输出通道及每轮两拍重构。每ADC复用8个输出列；每次只求值16列，再换组重新建立，无跨组模拟保持。每输入bit遍历32个行组和8个输出组，故共享API给出
\[
 N_E=N_C=N_{\rm rec}=8\times32\times8=2048,\qquad
 N_{\rm dig}=4096,\qquad N_{\rm scalar}=262144.
\]
这里的标量转换次数不是逻辑输入字节，整个过程仍只服务256 Byte。

$V_{\rm MAC}$ 随状态组合变化，\textbf{SAR的原始电压码不能直接充当 $0\ldots8$ 的计数}。每ADC通道配置9等级的标定阈值或小型译码表，将电压码映为近似部分和；需覆盖所选列的偏置/增益差异。第1拍并行完成128个等级译码和16组八平面位权合并，第2拍按输入bit权重累加到对应输出寄存器。这个有限等级译码与加法满足所选 $T_D$ 是参考实现的资源/时序条件；原四阶段VSR-VSA由该前端和SAR替换，不再额外累加。SLC电导离散、泄漏、状态组合非线性及相邻电压分布重叠仍形成近似误差，二状态终态通过不能消除它们。

\sid{PCM-03} 报告完整access范围3.25--15.9 ns，在0.95 V下8b IN/8b W/19b OUT的Shmoo点为14.3 ns（PDF p.1、p.2 Fig.11.3.6）。它已包含原宏前端、感测和局部逻辑；不能从中拆出一个实测cell延迟，也不能把14.3 ns再完整叠到公共输入、ADC和重构上。这里仅用它支持八项电压前端的数ns到数十ns量级，选取
\[
 F_A=T_I+t_{m,A}=10/20/50\ns
\]
作为\emph{含输入形成和公共复位}的总前端预算，覆盖WL/列选通、固定电流建立及电压节点稳定；原文没有单独测出这三个数。短端要求较快前端，参考和长端给电压建立及接口负载较多预算。公共SAR再用 $T_A=10/20/50$ ns，两拍数字重构用 $T_D=2/5/10$ ns。复算程序调用共享A模板时填 $t_{m,A}=F_A-T_I$，避免重复计算输入形成。

全部阶段顺序执行；输入捕获与128个结果容器最终提交合计两拍。因此
\begin{equation}
 \boxed{\Delta_S=2048(F_A+T_A+2T_D)+2T_D,\qquad
 \rho\approx256\,\mathrm{Byte}/\Delta_S.}
\end{equation}
参考点为 $2048(20+20+10)+10=102410$ ns。前端与ADC各占40960 ns，数字译码/重构20480 ns，边界10 ns。32个八行组覆盖原生256输入，八个输出组与位串行共同决定读占用；它不表示PCM单cell读需数十微秒。八行保留电压等级和信号裕量的证据，不是材料活动行数上限。

\section{二状态写序列、终点与跨实现桥接}

\sid{PCM-01} Methods（PDF p.7）给出RESET125 ns、700 $\mu$A，以及SET250 ns、尾沿50 ns、125 $\mu$A。其措辞没有明确250 ns是否已含50 ns尾沿，本文保守计SET总波形300 ns，另记录总250 ns的敏感性。每物理批先RESET所有目标cell，再只对目标bit为1的cell施加SET；SET掩码之外保持RESET，固定调度不因数据稀疏而缩短。每批RESET、SET各一次，随后实际读取全部cell的最终状态。

二状态的依据是原文明确分离的SET/RESET分布：绝大多数核心的unit-cell可达 $|G_{\rm RESET}|<5$、$G_{\rm SET}>50$ ADC counts，报告超过99\%的unit-cell满足该标准，一个核心为98.4\%（PDF p.4、p.14 Fig.2a）。这支持大幅分离的端态，而不证明本参考设计256个cell全部一次成功，更不证明相同读出电流。原4-PCM unit-cell和CCO码值不能作为本设计单PCM/SAR的通用量纲。实现时在同一读偏置和前端下标定HRS上阈值、SET下阈值及量程检查，保留互不相交的状态窗；失败必须返回状态，不能计为成功的 $B_R$。

主情景描述的是\textbf{经过器件表征的正常单次端态更新服务}，不估造未报告的批良率、额外重试分布或最大完成时间。其有限预算有具体依据：原文对端态确实使用一次完整SET/RESET波形，而连续中间电导才进一步采用125 ns迭代脉冲、每次read/verify、5 ADC counts误差窗及最多30次迭代。把30次直接搬到二进制程序没有依据。这里完整保留两种端态波形与真实终验；若目标实现经常需要追加脉冲，就必须按实际批次数增加写占用，当前 $\tau$ 不能当作该失败尾部的持续带宽保证。

\sid{PCM-06} 的60/120/240 ns脉宽实验显示器件间、周期间变异，且中间状态变异较大（PDF pp.3--4）；其512列行并行使用公共幅度DAC和每列独立脉宽，连续模拟闭环需要FPGA计算下一轮。\sid{PCM-05} 的独立Ge-rich GST测量采用75 ns脉冲，但材料、幅值和选择器条件不同（PDF p.3）。它们支持几十到几百ns的编程波形量级，并不把60 ns或75 ns变成本例完整周期。1000次是实验统计序列，1.2 ns/tick是输出编码；二者均未进入每更新必需操作数。

每批专用驱动额外计三段 $g=10/20/50$ ns：IDAC/路由建立、RESET后的冷却及SET偏置转换、SET完整尾沿结束后的路径隔离与低压回读偏置恢复。这是波形之外的\emph{明确工程预算}，不是原文测得的热恢复常数；其量级相对125--300 ns波形较小，参考点三段共60 ns。其中SET的50 ns下降尾沿已经包含在300 ns内，第三段从尾沿完成之后开始，不再支付同一下降波形；后续每次二状态校验重新执行完整本地前端建立。没有复制RRAM的微秒级高压余量，也没有将材料亚ns切换或FPGA等待补成周期。

\section{32-Byte 行内分组与复用 ADC 的完整读验}

固定输入行 $r\in[0,255]$ 和列内槽 $s\in[0,3]$，事务的32个逻辑权重为
\[
 (r,c),\qquad c=4h+s,\quad h=0,\ldots,31.
\]
32个IDAC分别负责一个4列组，每次仅一个平面活动；顺序遍历八平面，每批32个cell，共256个cell，完成32个INT8权重，因此 $B_R=32$ Byte。两拍128-bit数据装入256-bit缓冲，首拍与命令合并，准备/完成及装入总占用 $T_{\rm front}=3T_D$，不是由接口宽度推出32-cell并行。

校验复用当前平面的16 ADC。ADC号 $a\in[0,15]$ 负责8列组，两批依次选择 $c=8a+4k+s$、$k=0,1$，同时选中同一WL行。每批16个ADC各读取一cell；第二批\emph{重新复位、建立并读出}，不保留前一批的模拟结果，不假设32个免费hold槽。原生WL与列偏置在每个活动列只选一cell，此时不启用CIM的八WL并行模式。

单cell二状态终验采用\textbf{同一 $I_M$、钳位和列电压通路}，每批完整重复 $F_A$，随后公共SAR转换 $T_A$ 和一拍16路终态比较。这个桥接不是因为“binary必然更快”，而是八行SLC求值本身已经包括 $q=1$（一个SET、其余HRS）和 $q=0$（全HRS/off）两端：所选 $F_A$ 必须覆盖这些最弱有效等级的建立和量程，不能只对八个SET的强信号成立。同row校验时同一访问管的gate驱动电平与计算模式相同，最多一cell导通；列负载、$I_M$ 和钳位保持，未选支路泄漏/偏置的变化由校验专用两阈值标定处理。无导通/高阻端由钳位提供确定的低端或高端电压，而不等待未定义的浮动节点。上述兼容性与最弱码建立是本参考前端的工程条件；若某实现单cell需要更长建立，应修改 $F_{V}$，而不是沿用主值。

保留完整 $F_V=F_A$ 没有把原14.3 ns整体access重复加入，也没有假定全部八cell总电流始终更大；判决只接受两个分离端态，不要求模拟中间电导精调。每批重新建立确保不依赖保持槽或前批结果。令 $t_V=F_A+T_A+T_D$，完整resident服务为
\begin{equation}
\boxed{\begin{aligned}
 \Delta_R&=3T_D+8\left[T_D+3g+125\ns+300\ns+2t_V\right],\\
 \tau&\approx32\,\mathrm{Byte}/\Delta_R,\qquad
 \RI^*\approx8\Delta_R/\Delta_S.
\end{aligned}}
\end{equation}
内层一拍用于批地址/模式配置，三段专用转换及两种端态波形随后执行，再依次完成两批终验。八个权重平面、RESET/SET、16次校验均只增加服务时间，逻辑写payload只计一次。RESET后不作独立中间校验，因为最终校验已经检查目标0和目标1，程序没有以RESET中间结果决定下一脉冲。

逐一遍历256行和4个列内槽，得到1024个32-Byte事务；每个 $(r,c)$ 唯一对应行 $r$ 与槽 $s=c\bmod4$。完整矩阵无遗漏、无重复，总payload为 $1024\times32=32768$ Byte，总时间 $T_R=1024\Delta_R$。脚本检查代表性 $r,s,w$ 下，只有已选WL与32个有编程压差的列相交才得到目标；同一行半选cell无压差，其他行gate关闭。这是逻辑选通检查，回流/电流合规仍依赖上述实际外围。

任意连续或稀疏小更新必须按实际列组利用率重算，不能自动套用32 Byte满利用率。增加写头须同时扩驱动、路由和回流，不作为本固定配置快慢端。RESET已逐批支付，无额外块擦摊销。

\section{成对结果、主导因素与边界}

% Generated by scripts/check_pcm.py --emit.
\begin{table}[htbp]\centering\small
\begin{tabular}{@{}lrrrrr@{}}\toprule
情景 & $\Delta_S$ ($\mu$s) & $T_R$ (ms) & $\rho$ (MB/s) & $\tau$ (MB/s) & $\RI^*$\\\midrule
短 & 49.156 & 4.110 & 5.208 & 7.972 & 0.6533\\
参考 & 102.410 & 4.767 & 2.500 & 6.874 & 0.3636\\
长 & 245.780 & 6.625 & 1.042 & 4.946 & 0.2106\\
\bottomrule\end{tabular}
\caption{原生 $K=256,N=128$、八行电压式近似 ACIM 的配对服务；$T_R$ 是32768 Byte完整矩阵装载，局部事务为32 Byte。MB 为十进制。}\label{07_pcm:tab:results}\end{table}


参考点手算为
\[
 \Delta_R=15+8[5+60+125+300+2(20+20+5)]=4655\ns,
\]
\[
 \rho=(256/102410)\times10^3=2.4998\,\mathrm{MB/s},\qquad
 \tau=(32/4655)\times10^3=6.8743\,\mathrm{MB/s},
\]
\[
 \RI^*=(256/32)(4655/102410)=0.363636.
\]
完整矩阵装载 $T_R=4.766720$ ms，$\tau=32768/T_R$；两表接口 $U^*=T_R/\Delta_S=128\RI^*\approx46.55$。每16 KiB平均更新成本仅是 $T_R/2$，不是实际小请求延迟。主范围是固定模式下三个成对预算，不能将独立最短读与最长写随意组合。共享 $T_D/T_A$ 与由其推导的控制时间必须联动，未将不兼容读写端点交叉组合成物理范围。成对ridge变化来自固定编程波形与读前端/外围占用的不同比例。较低ridge也包含八行电压组织使streaming重复2048轮、降低 $\rho$ 的后果；它不表示PCM适合所有动态更新。是否适用还取决于绝对两路能力、32权重行内列组利用率与完整矩阵重载成本。

\begin{table}[htbp]\centering\small
\begin{tabular}{@{}lrr@{}}\toprule
参考点 resident 占用 & ns & 占比 (\%)\\\midrule
命令、数据装入及提交 & 15 & 0.32\\
八批选通/模式控制 & 40 & 0.86\\
三段专用驱动转换 & 480 & 10.31\\
RESET 脉冲 & 1000 & 21.48\\
SET 平台及尾沿 & 2400 & 51.56\\
二状态重新建立/读出 & 320 & 6.87\\
替代 SAR 转换 & 320 & 6.87\\
终态比较 & 80 & 1.72\\
\bottomrule\end{tabular}
\caption{完整写服务分解。接口与内部读验不增加逻辑写入字节。}\label{07_pcm:tab:write}\end{table}


参考点的两种编程波形共3400 ns，占写服务73.0\%；三段驱动转换共480 ns，占10.3\%；16次二状态前端建立与SAR各320 ns，合计13.7\%。因此此二状态主程序由SET/RESET波形主导，前端和公共外围也已完整支付。三档 $g$ 与外围快慢仅为成对选取，没有必须相等的物理关系。固定参考读前端和公共外围，仅把每段工程驱动余量 $g=20$ ns 改为10/50 ns，得到 $\Delta_R=4.415/5.375\us$、$\tau=7.248/5.953$ MB/s、ridge约0.345/0.420；$\rho$ 不变。该独立对照量化工程预留的影响，不把 $g$ 称作PCM不可消除的材料时间。SET总长若按250 ns解释，$\Delta_R$降到4.255\,$\mu$s，$\tau$约7.521 MB/s，ridge约0.332；这约8.6\%差异保留为波形解释的不确定性。

唯一组织对照保持相同状态、32-IDAC、16-ADC/平面、每次32-Byte事务及所有program步骤，把每次终验前端由 $F_A$ 改成\textbf{保守模拟弱信号读验}的768 ns。其依据是\sid{PCM-01} PDF p.8的小电导校验：0.2 V、512 ns读脉冲和256 ns预充电，原CCO在脉冲内编码。本对照将768 ns整体作为新弱信号前端配额，其后用共享SAR替换原CCO；没有另计原CCO转换，也不声称原768 ns是纯器件延迟或SLC必需时长。

参考对照 $\Delta_R=16.623\us$、$\tau=1.925$ MB/s、ridge=1.298，streaming不变；比本地主二状态路径多 $16(768-20)=11968$ ns，写服务约慢3.57倍。该对照中的前端读验占73.9\%，说明把精细模拟小电导组织强加给粗端态，会改变写侧主导项。若所选前端不能在相同偏置/钳位下满足 $q=0/1$ 的建立条件，弱信号方案是可追溯的保守预算；两者差别不能被解释成PCM材料固有写速度范围。补充任何更长单cell建立条件都须同时更新主读的最弱码兼容性判断。

\needspace{5\baselineskip}
读写使用同一状态和ADC，本文分别估计两路服务能力，不保证同时实现各自峰值。PCM读后无需例行restore；本地服务窗口没有加入周期刷新。保持漂移不会因选择二状态而完全消失，长期再标定或重编程如有需要，应额外计入占用，不能增加workload的逻辑更新字节；耐久与寿命另列可行性约束。

R1保留原生 $K=256,N=128$、共同逻辑字节、近似ACIM合同和单域；R2明确八行组来自代表性电压实现；R3保持128 ADC、16数字通道和两拍重构，增加必要9等级标定译码；R4令前端预算含 $T_I$ 并替换原VSR/CCO，不重复叠完整原宏；R5以实际32-IDAC、256-bit缓冲、行内列组事务及读验mux计算写入。最大不确定性是端态波形跨实现适配，以及同一电压前端对单cell $q=0/1$ 的建立兼容性；其条件和弱信号组织对照均留在输入JSON及上述预算内。脚本导入共享API，检查实际映射、两批重新读验、阶段覆盖、单位、成对结果、源PDF及共享文件哈希；默认只读，显式\sid{--emit}才更新派生数据和数字表。
